% Section 7 — The decisive experiment, specified in advance
\section{The decisive experiment, specified in advance}
\label{sec:protocol}

This section specifies, before any of it is run, the experiment that would answer the question of
\cref{sec:intro}. Its design follows one principle: \emph{every way the result could be explained
without learning from experience gets its own control arm.}

\subsection{Unit, endpoint and power}

\paragraph{Unit} The unit of analysis is the dataset \emph{family}: a dataset together with its
versions, aliases and derived targets. Families are drawn from the sealed holdout, which no model,
teacher or catalog in this work has touched, and are frozen, with their identities hashed, before
any arm runs. Each family contributes $T_f$ tasks evaluated over $R=9$ outer splits (three split
seeds, three folds).

\paragraph{Utility} For a committed program $a$ on an outer split, let $e$ be its held-out RMSE and
$b$ that of the training-mean predictor, with $\varepsilon$ computed from training targets only:
\begin{equation}
u = \begin{cases} \mathrm{clip}\bigl(\log\frac{\max(b,\varepsilon)}{\max(e,\varepsilon)},\,-1,\,1\bigr) & \text{valid program,}\\
-1 & \text{any failure.}\end{cases}
\label{eq:utility}
\end{equation}
A failure is any refusal, compile or runtime error, timeout, memory exhaustion, non-finite output or
protocol violation. Failures are never dropped; infrastructure retries are distinct from method
failures and are logged.

\paragraph{Primary endpoint} With $U_{f}$ the family-mean utility,
$\Delta = \tfrac1F\sum_f (U_f^{\mathrm{IZ}} - U_f^{\mathrm{B}})$, where B is the strongest
selection or retrieval arm (\cref{tab:arms}), chosen on validation families and named before
unblinding. The primary hypothesis is superiority with margin $m=0.02$ utility units: the lower bound
of the two-sided 95\% family-bootstrap interval (10,000 resamples, BCa \citep{efron1987bca}) must
exceed 0.02. The margin is not a 2\% RMSE improvement; the raw log-RMSE ratio is reported
separately. Elo, where shown, is refitted inside every bootstrap resample and is never the primary
statistic.

\paragraph{Power} Planning for a true effect of 0.05 with a family-level standard deviation of 0.15
requires $\lceil (1.96+0.84)^2\, 0.15^2/(0.05-0.02)^2\rceil = 196$ families; we register a floor of
240 and a ceiling of 640, with 60 further families reserved for mechanism tests. The standard
deviation is an assumption, not a pilot measurement: at 0.25 the same calculation asks for 545
families, so the roster is fixed by development-only simulation before any confirmation outcome is
opened, and a study that cannot reach adequate power is declared exploratory in advance. The
estimand, interval method, multiplicity rules and resource envelope are specified in full in
Supplement~\ref{S-sup:statistics}. In a ten-dataset
pilot like \cref{sec:pilot}, a sign test reaches significance only when one method wins at least
nine of ten datasets, which is why it is labelled a pilot.

\paragraph{Commit first, evaluate independently} For each outer-training context the proposer
receives only permitted training information and a common resource card, and commits its source,
constants, dependencies and component identities before any candidate-specific execution. A
deterministic structural check may refuse an output, but a refusal consumes the attempt: there is no
resampling and no replacement. Compilation follows commitment, so a compiler failure remains a
failure of the primary attempt, and diagnostics never return to the proposer. An isolated supervisor
records the order of events and the source hashes rather than trusting a proposer's account of what
it did; the evaluator then compiles, fits and scores the unchanged source, and only the scorer holds
the test labels. Selection among alternatives inside \code{fit} is disclosed as fitting-time search
rather than hidden inside the proposal budget.

\paragraph{Matched information} Every conditional arm --- the learned proposer, the frontier models
asked for one program per family, the retrieval and selection controls --- receives the same
deterministic, training-only view of each task: sampled rows, neutral feature identifiers, permitted
targets and the resource card, hashed before use. Size limits and any label-independent feature
reduction are fixed in advance, and no provider may truncate its input silently. A separate arm can
then measure what that compression costs. The committed programs themselves fit on the whole
outer-training partition.

\subsection{Arms}

\begin{table}[t]
\caption{Arms of the decisive experiment, all at $K=1$ on the same untouched families and the same
resource card. A5 is the control that distinguishes composition from a portfolio.}
\label{tab:arms}
\centering\sffamily\footnotesize\lining
\begin{tabularx}{\linewidth}{@{}llL@{}}
\toprule
Arm & Track & What it isolates \\
\midrule
A1 InnovationZero & C & open-support proposer trained on $J_1$ (\cref{eq:j1}); the treatment \\
A2 Catalog-bound policy & C & selection only (\cref{prop:closed}); same encoder and training budget \\
A3 Retrieval & C & nearest training family in encoding space, then its incumbent program \\
A4 Best single program & U & the one kernel with the best mean utility over training families \\
A5 Distilled bank & C & A1's own programs, collected on a separate development pool, frozen as a bank, plus a selector trained with A1's encoder and budget \\
A6 Shuffled context & C & A1 given the encoding of a different dataset: does it read the data? \\
A7 Frontier entrants & U & pristine re-issue of the challenge (E0), three attempts per model, all reported \\
A8 Frontier, conditional & C & frontier models given the information-matched data packet, one program per family \\
A9 Search at budget $k$ & S & the engine's proposer with $k\in\{1,4,16,64,256\}$ evaluations, with and without A1 as a seed \\
A10 Libraries & S & XGBoost, LightGBM, CatBoost, random forest at defaults and tuned to A9's wall clock \\
\bottomrule
\end{tabularx}
\end{table}

A5 is the decisive control. If a selector over A1's own frozen outputs matches A1 on untouched
families, A1's advantage is a good portfolio, not composition. If A1 wins, and wins most on families
whose best programs combine operations in ways absent from training (a compositional hold-out), that
is evidence of a capacity to construct. A6 tests whether the encoding is used at all. A8 receives,
in text, exactly the table A1's encoder reads --- at most 4,096 training rows and 256 features,
capped at 64 KiB --- so that any difference between them is not a difference in information.

\subsection{The InnovationZero curve}

The amortization question has a direct measurement. Let $U_{\mathrm{S}}(k)$ be the mean utility of
the search-assisted proposer with $k$ evaluations and $U_{\mathrm{IZ}}$ that of A1's single
proposal. The \emph{equivalent search budget} of the learned proposer is
\begin{equation}
k_{\mathrm{eq}} = \min\bigl\{ k : \mathbb{E}[U_{\mathrm{S}}(k)] \ge \mathbb{E}[U_{\mathrm{IZ}}] \bigr\},
\label{eq:keq}
\end{equation}
estimated with a joint family bootstrap and reported as ``not reached within 256 evaluations'' rather
than extrapolated. $k_{\mathrm{eq}}$ is the number of experiments a first proposal is worth. It is
the analogue, for algorithms, of the gap between AlphaGo Zero's raw network and its search; repeated
at several fractions of training history (none, a quarter, a half, all), it measures how much of
that worth comes from experience. \Cref{fig:curve} shows the design with the only point measured
today.

\begin{figure}[t]
\centering
% The InnovationZero curve as specified. Only the k = 256 engine point and the horizontal
% references are measured (pilot pool, evidence/pilot-2026-10-22/elo-deep.csv). Every other
% budget is pending (E2); no value is drawn for it.
\begin{tikzpicture}
\begin{axis}[iz,
  width=0.9\linewidth, height=0.5\linewidth,
  xmin=-0.5, xmax=8.6, ymin=1280, ymax=1820,
  xtick={0,1,2,4,6,8}, xticklabels={0,1,4,16,64,256},
  xlabel={Evaluation budget $k$ before the program is frozen (candidates executed on training rows)},
  ylabel={Elo (pilot pool)}, ymajorgrids,
  axis y line*=left, axis x line*=bottom, clip=false]
% frontier band and references
\addplot[draw=none, fill=izwarm, fill opacity=0.10, forget plot] coordinates {(-0.5,1506.2) (8.6,1506.2) (8.6,1761.4) (-0.5,1761.4)} \closedcycle;
\addplot[izwarm, line width=0.7pt, forget plot] coordinates {(-0.5,1761.4) (8.6,1761.4)};
\addplot[izwarm!80, line width=0.5pt, dashed, forget plot] coordinates {(-0.5,1649.0) (8.6,1649.0)};
\addplot[izwarm!60, line width=0.5pt, dashed, forget plot] coordinates {(-0.5,1506.2) (8.6,1506.2)};
\addplot[izgreen, line width=0.6pt, densely dotted, forget plot] coordinates {(-0.5,1482.2) (8.6,1482.2)};
\addplot[izgreen, line width=0.6pt, densely dotted, forget plot] coordinates {(-0.5,1326.0) (8.6,1326.0)};
\node[font=\sffamily\scriptsize\color{izwarm}, anchor=west] at (axis cs:8.7,1761.4) {Opus};
\node[font=\sffamily\scriptsize\color{izwarm}, anchor=west] at (axis cs:8.7,1649.0) {Sol};
\node[font=\sffamily\scriptsize\color{izwarm}, anchor=west] at (axis cs:8.7,1506.2) {Astra};
\node[font=\sffamily\scriptsize\color{izgreen}, anchor=west] at (axis cs:8.7,1482.2) {CatBoost};
\node[font=\sffamily\scriptsize\color{izgreen}, anchor=west] at (axis cs:8.7,1326.0) {XGBoost};
% pending budgets
\foreach \x in {0,1,2,4,6}{
  \edef\tmp{\noexpand\draw[izrule, line width=0.5pt, dash pattern=on 1.5pt off 1.5pt] (axis cs:\x,1300) -- (axis cs:\x,1800);
            \noexpand\node[font=\noexpand\sffamily\noexpand\tiny\noexpand\color{izgrey}, fill=white, inner sep=1pt] at (axis cs:\x,1300) {pending};}
  \tmp}
\node[font=\sffamily\scriptsize\color{izaccent}, align=left, anchor=south west, inner sep=1.5pt] at (axis cs:0.08,1765) {$k=0$: the zero-shot claim};
% the one measured engine point
\addplot[only marks, mark=*, mark size=3pt, izaccent] coordinates {(8,1630.9)};
\node[font=\sffamily\scriptsize\color{izaccent}, anchor=north east] at (axis cs:7.9,1622) {engine, measured};
\end{axis}
\end{tikzpicture}

\caption{The InnovationZero curve as specified. The one measured point is the engine at
$k=256$ on the pilot panel; the horizontal lines are the frontier entrants and two libraries on the
same panel. Every other budget, and the $k=0$ point that is the zero-shot claim, is pending; nothing
is drawn for them \pend.}
\label{fig:curve}
\end{figure}

\subsection{Novelty and Move 42}

\paragraph{A novelty ladder} ``New'' is graded, and each grade licenses specific words
(\cref{tab:novelty}). Before unblinding, the complete reference archive --- teachers, catalog,
retrieval bank, decomposition corpus and every program used in training --- is frozen with its
canonicalizer and hashes. A program is \emph{out of archive} only if no archived program matches it
after canonicalization, and the difference is more than names, constants, fitted state, dead code or
emitter expansion; two blinded reviewers and an adjudicator must agree on what reusable structure
changed. The \emph{useful out-of-archive rate} $p_{\mathrm{UN}}$ is the fraction of all attempts,
failures included, that are valid, out of archive and useful (no more than 0.95 times ridge
regression's error, and within 5\% of the best arm's).

\begin{table}[t]
\caption{The novelty ladder: what each level requires and what it permits us to say.}
\label{tab:novelty}
\centering\sffamily\footnotesize\lining
\begin{tabularx}{\linewidth}{@{}lLL@{}}
\toprule
Level & Requirement & Permitted wording \\
\midrule
N0 & a new source hash & ``new bytes'' \\
N1 & canonical structure differs & ``a different program'' \\
N2 & outside the complete frozen archive & ``not retrieval'' \\
N3 & composes relations held out of training & ``compositional transfer'' \\
N4 & behaves differently on independent probe tasks & ``a functionally new program'' \\
N5 & removing or replacing its mechanism lowers held-out utility & ``a useful new mechanism'' \\
N6 & an expert prior-art audit finds no precedent within stated limits & ``contextually novel'' \\
\bottomrule
\end{tabularx}
\end{table}

\begin{definition}{Move 42}
\label{def:move42}
A \emph{Move 42} is a program committed under the track-C contract that (i) lies outside the frozen
archive (N2), (ii) improves held-out utility over the strongest arm on its family, and (iii) contains
a mechanism whose removal, or replacement by the nearest known component at matched capacity, lowers
utility on a separately reserved set of 60 families, with the lower bound of the 95\% interval above
zero and the effect not confined to one data source. The case is selected by a rule fixed in
advance --- the lowest committed source hash among qualifying programs --- not by inspection.
\end{definition}

The name marks a debt and a distinction. Move~37 was a move in a game whose rules people wrote, found
by search at the moment of play, and judged surprising because an imitation model gave it one chance
in ten thousand. Move~42 is an answer the data gives about itself --- an algorithm, not a move --- found
without executing anything, and judged by whether it works and why. Rarity under a frontier model's
prior can be reported as a secondary statistic: zero occurrences in $n$ independent draws bound the
probability by roughly $3/n$, so a one-in-ten-thousand claim needs about 30,000 draws, and even then
rarity is not novelty. If no program meets \cref{def:move42}, we will report that no Move~42 was
found.

\subsection{Claim ladder and stop rules}

Results are reported against a ladder fixed in advance (\cref{tab:ladder}); each rung licenses one
claim, a rung is passed only if every rung below it is, and missing evidence is never averaged into
success. Three stop rules are fixed with it: if A6 matches A1, we do not claim data-conditioned
proposal; if A1's programs remain inside the catalog, we stop claiming open synthesis; and if A1's
utility improves only after any selection among complete candidates, it is reclassified as track S.

\begin{table}[t]
\caption{The claim ladder. Today only M0 is partially met.}
\label{tab:ladder}
\centering\sffamily\footnotesize\lining
\begin{tabularx}{\linewidth}{@{}llLl@{}}
\toprule
Rung & Gate & Requirement & Claim permitted \\
\midrule
M0 & custody & every attempt reconciled, every artifact hashed & an auditable instrument \\
M1 & boundary & $K=1$, $E=0$, $R=0$ recorded by the supervisor & a zero-feedback proposer \\
M2 & completion & lower bound of executable rate $\ge 95\%$ & a reliable proposer \\
M3 & utility & $\Delta$ lower bound $>0.02$ against B & a useful first proposal \\
M3b & attribution & A1 $>$ A5 and A1 $>$ A6 & learned from experience, not a portfolio \\
M4 & novelty & lower bound of $p_{\mathrm{UN}} > 10\%$ & useful synthesis beyond the archive \\
M5 & mechanism & \cref{def:move42} met and replicated & a Move~42 \\
M6 & economics & lower cost per useful algorithm at matched quality, fully amortized & an efficiency claim \\
M7 & self-improvement & successive generations beat matched static controls on fresh families & bounded self-improvement \\
\bottomrule
\end{tabularx}
\end{table}

\paragraph{Resources} The planning envelope is one CPU and 8~GiB for preprocessing, compilation,
fitting and prediction, with limits of 120, 120, 600 and 60 seconds; encoding may use one disclosed
GPU of at most 24~GiB for 120 seconds, and proposal through commitment has a 600-second ceiling and
one trajectory. At 240 families, three split seeds and three folds give 2,160 contexts, and so 2,160
fits, per arm: at most 360 CPU-hours each at the fitting ceiling, or 3,600 for the ten arms of
\cref{tab:arms} before A9's additional budgets. That is a bound, not a forecast.
Every scheduled attempt must reconcile to a commitment or a typed failure, and no timeout, missing
receipt or dropped difficult task may leave the denominator.

\paragraph{Before the decisive experiment} Two cheaper experiments come first. \textbf{E0} re-issues
the challenge to frontier models in fresh sessions with only the participant packet, metered usage
and three attempts per model, replacing every estimate in \cref{sec:yield}. \textbf{E10} decomposes
the E0 inventions into typed operations, adds them to the kernel language and lets A1 learn when to
use them; it tests whether general inventions and per-dataset composition are better together than
apart (\cref{fig:routes}).
